% Chapter 4

% \doublespacing
% \setlength{\parindent}{0.5in}
% \setlength{\parskip}{4pt}
% \justifying

\chapter{\uppercase{SYSTEM DESIGN AND IMPLEMENTATION}}

\section{\uppercase{ARCHITECTURAL OVERVIEW}}

The proposed zero-day threat detection system is designed as a modular pipeline that transforms raw system events into behavioral insights through graph-based anomaly detection. The architecture prioritizes operational efficiency by maintaining bounded computational and memory costs while processing heterogeneous event streams in real time. Key architectural principles include user-space operation (eliminating kernel modifications), multi-threaded event collection (enabling concurrent monitoring), graph-structured representation (capturing behavioral relationships), and hybrid anomaly detection (combining reconstruction error with structural analysis). The system encompasses four primary processing stages: real-time event collection, behavior graph construction, graph autoencoder modeling, and hybrid anomaly detection with statistical threshold-based alerting.

This design enables the system to detect previously unknown attack patterns by identifying behavioral deviations from established baselines, rather than relying on signature-based detection or historical attack databases. The modular architecture permits independent development, testing, and optimization of each component while maintaining strict integration contracts between stages.

\section{\uppercase{SYSTEM DESIGN METHODOLOGY}}

The architecture of the proposed system follows a modular design consisting of several interconnected components that together form the intrusion detection pipeline. The system begins with a real-time system event collection module that continuously monitors host activities such as process execution, file system interactions, and network communications. These events are captured and stored in a temporary buffer while preserving their temporal order.

The collected event stream is then processed by the behavior graph construction module. This module transforms raw system events into structured graph representations where nodes represent system entities such as processes, files, or network sockets, and edges represent interactions between these entities. This graph-based representation enables the system to capture complex behavioral relationships that may indicate malicious activity.

Once the behavior graphs are generated, they are analyzed using a graph autoencoder-based anomaly detection model. The graph autoencoder learns patterns of normal system behavior by reconstructing input graphs and identifying deviations through reconstruction error. Graphs that cannot be accurately reconstructed are treated as potential anomalies.

Finally, the system includes a visualization and alert generation component that presents detection results to the analyst. This module displays detected anomalies, system metrics, and behavioral graphs through a web-based dashboard, enabling users to monitor system security in real time and respond to suspicious activities effectively.


\section{\uppercase{REAL-TIME DATA COLLECTION}}

The first stage of the system is responsible for collecting system events in real time. The collector monitors process creation, file system operations, and network connections to obtain a comprehensive view of host behavior. Each event is timestamped and stored within a bounded queue to preserve temporal ordering.

The event collector uses multiple monitoring threads to capture different types of system activities simultaneously. This design ensures minimal latency while maintaining continuous event acquisition.

The procedure followed by the real-time system event collector is described in Algorithm 4.1. The algorithm outlines how different system activities are monitored and stored in a temporary event buffer for further analysis.

\begin{algorithm}
\caption{Real-Time System Event Collection}
\begin{algorithmic}[1]

\State Set polling interval $\delta = 100$ ms
\State Initialize event buffer with capacity $B = 10,000$

\While{system is active}

\State Monitor newly created processes
\State Monitor network connection events
\State Monitor file system operations

\For{each detected event}
\State Record event details
\State Insert event into buffer
\EndFor

\State Wait for $\delta$ milliseconds

\EndWhile

\State Sort events by timestamp
\State Return event stream

\end{algorithmic}
\end{algorithm}

The collected events form the raw behavioral data used for constructing interaction graphs in the next stage of the system. The event buffer operates with defined capacity $B = 10,000$ and polling interval $\delta = 100$ milliseconds, ensuring that system resources remain predictable and detection latency remains bounded even under high event volume.



\section{\uppercase{BEHAVIOR GRAPH CONSTRUCTION}}

The event stream collected from the host system is transformed into behavior graphs. In these graphs, nodes represent system entities such as processes, files, and network sockets, while edges represent interactions between them.

Events are grouped into fixed time windows to capture short-term system activity. For each time window, a graph is constructed that represents the relationships between system entities.

The transformation of raw system events into structured behavior graphs is a critical step in the proposed intrusion detection framework. This stage captures the relationships between processes and system resources by representing them as nodes and edges in a graph structure. During graph construction, node features are computed to characterize entity behavior, including temporal statistics (event frequency, inter-arrival times), structural properties (degree values, connectivity patterns), and entity type indicators.

The detailed procedure used to construct these graphs from the collected event stream is presented in Algorithm 4.2.

\begin{algorithm}
\caption{Behavior Graph Construction}
\begin{algorithmic}[1]

\State Input event stream $E$
\State Define time window $W = 5$ seconds

\State Partition events into time windows

\For{each window}

\State Initialize empty graph $G$

\For{each event}

\State Identify process node
\State Identify resource node
\State Add nodes to graph
\State Create edge between nodes

\EndFor

\For{each node}

\State Compute node type feature
\State Compute normalized degree values
\State Compute temporal activity rate

\EndFor

\State Store graph $G$

\EndFor

\State Return graph sequence

\end{algorithmic}
\end{algorithm}

The generated graphs preserve structural relationships between processes and system resources. These graphs serve as input to the graph autoencoder model used in the next stage of the system.

Graph-based representations enable the detection model to capture complex relationships between system entities. Among the 14 structural features extracted, edge density is computed as shown in Equation \ref{eq:edge_density}:

\begin{equation}
\rho = \frac{2|E|}{|V|(|V|-1)} \quad \text{(Equation 4.1)}
\label{eq:edge_density}
\end{equation}

Equation 4.1 denotes the normalized edge density metric, which measures the ratio of actual edges to maximum possible edges in the graph, capturing the connectivity concentration of process-resource interactions.

Furthermore, dividing the event stream into fixed time windows ensures that the system analyzes localized behavioral patterns rather than a continuous unstructured stream. This improves computational efficiency while maintaining meaningful temporal context in the behavior graphs.



\section{\uppercase{GRAPH AUTOENCODER MODEL}}

To detect anomalous system behavior, the proposed system employs a Graph Autoencoder. The model learns the structural patterns of benign system activity by reconstructing the input graph.

The encoder compresses node features into a lower-dimensional embedding space, while the decoder reconstructs the graph adjacency structure from learned representations. The autoencoder is trained exclusively on benign system graphs to establish a baseline model of normal behavior. The model minimizes the reconstruction loss function as given in Equation \ref{eq:loss}:

\begin{equation}
\mathcal{L} = \|A - A'\|_F^2 \quad \text{(Equation 4.2)}
\label{eq:loss}
\end{equation}

Equation 4.2 denotes the Frobenius norm of the difference between the original adjacency matrix and the reconstructed matrix, quantifying how well the autoencoder learns normal behavior patterns. where $A$ is the original adjacency matrix, $A'$ is the reconstructed adjacency matrix from the decoder output, and $\|\cdot\|_F$ denotes the Frobenius norm.

The training and detection process of the graph autoencoder is summarized in Algorithm 4.3.

\begin{algorithm}[H]
\caption{Graph Autoencoder Training and Detection}
\begin{algorithmic}[1]

\State Input behavior graphs
\State Initialize autoencoder model
\State Initialize optimizer

\For{each training epoch}

\For{each graph}

\State Extract node features
\State Normalize adjacency matrix
\State Compute node embeddings
\State Reconstruct adjacency matrix
\State Calculate reconstruction loss
\State Update model parameters

\EndFor

\EndFor

\State Compute anomaly score using reconstruction error

\end{algorithmic}
\end{algorithm}

Graphs that exhibit higher reconstruction errors are considered potential anomalies. The reconstruction error threshold is determined through statistical analysis of training data, allowing the model to adapt to the specific characteristics of the monitored system.


\section{\uppercase{HYBRID ANOMALY DETECTION}}

While graph autoencoders are effective at learning normal behavioral patterns, relying solely on reconstruction error may not always capture subtle structural anomalies. Therefore, the proposed system incorporates additional structural graph features into the detection process.

The hybrid detection mechanism is presented in Algorithm 4.4, where reconstruction-based anomaly scores are combined with structural graph metrics.

\begin{algorithm}[H]
\caption{Hybrid Anomaly Detection}
\begin{algorithmic}[1]

\State Input behavior graph $G$

\State Extract 14 structural features:
\State \ \ $F_0$-$F_2$: Node type distribution (process, file, socket counts)
\State \ \ $F_3$-$F_6$: In/out-degree statistics (average and maximum)
\State \ \ $F_7$-$F_9$: Node type ratios (normalized proportions)
\State \ \ $F_{10}$-$F_{12}$: Unique entity counts (processes, files, sockets)
\State \ \ $F_{13}$: Edge density ($e/\binom{n}{2}$)

\State Compute Z-scores for each feature:
\State $Z_i = \frac{F_i - \mu_i}{\sigma_i}$ for $i = 0, 1, \ldots, 13$

\State Compute structural anomaly score (Eq. \ref{eq:struct_score})

\State Compute reconstruction error (Eq. \ref{eq:recon_error})

\State Combine scores (Eq. \ref{eq:hybrid_score})

\If{$S_{hybrid}$ exceeds threshold $T$}

\State Flag graph as anomaly

\Else

\State Mark graph as normal

\EndIf

\State Return detection result

\end{algorithmic}
\end{algorithm}

The structural and reconstruction-based anomaly scores are computed according to Equations \ref{eq:struct_score}, \ref{eq:recon_error}, and \ref{eq:hybrid_score}:

\begin{equation}
S_{\text{struct}} = \frac{1}{14}\sum_{i=0}^{13} |Z_i| \quad \text{(Equation 4.3)}
\label{eq:struct_score}
\end{equation}

Equation 4.3 denotes the structural anomaly score computed as the mean absolute Z-score across all 14 structural graph features, quantifying deviation from normal graph structure patterns.

\begin{equation}
S_{\text{recon}} = \|G - G'\|_2 \quad \text{(Equation 4.4)}
\label{eq:recon_error}
\end{equation}

Equation 4.4 denotes the reconstruction error computed as the L2 norm between the original graph and its reconstruction by the autoencoder, measuring how anomalous a graph is based on learned normal patterns.

\begin{equation}
S_{\text{hybrid}} = 0.5 \cdot S_{\text{struct}} + 0.5 \cdot S_{\text{recon}} \quad \text{(Equation 4.5)}
\label{eq:hybrid_score}
\end{equation}

Equation 4.5 denotes the hybrid anomaly score combining structural and reconstruction-based detection with equal weighting, providing a balanced assessment of behavioral anomalies. where $Z_i$ represents the normalized Z-score for feature $i$, $G$ is the original behavior graph, $G'$ is the reconstructed graph from the autoencoder, and the weighting coefficients are set to 0.5 for equal contribution from both detection mechanisms. Graphs are flagged as anomalies when the hybrid score exceeds the statistically determined threshold.

Structural graph features capture 14 properties including node type distribution, in-degree and out-degree statistics, unique entity counts, and edge density. These characteristics provide comprehensive insight into system behavior that may not be fully captured by reconstruction error alone.

By combining reconstruction error with structural graph analysis as shown in Equation \ref{eq:hybrid_score}, the hybrid detection approach improves detection robustness and reduces false positives.


\section{\uppercase{THRESHOLD DETERMINATION AND TUNING}}

A critical aspect of anomaly detection is determining appropriate thresholds that balance sensitivity and specificity. The proposed system employs a statistical thresholding strategy based on analysis of the training data.

The threshold determination procedure is detailed in Algorithm 4.5.

\begin{algorithm}[H]
\caption{Statistical Threshold Determination}
\begin{algorithmic}[1]

\State Input benign graph set $G_{benign} = \{G_1, G_2, \ldots, G_n\}$
\State Initialize empty score list $S$

\For{each benign graph $G_i$}

\State Compute hybrid anomaly score $s_i$
\State Append $s_i$ to $S$

\EndFor

\State Compute mean $\mu = \frac{1}{n}\sum_{i=1}^{n} s_i$
\State Compute standard deviation $\sigma = \sqrt{\frac{1}{n}\sum_{i=1}^{n}(s_i - \mu)^2}$

\State Set detection threshold using Eq. \ref{eq:threshold}

\State Return threshold value $T$

\end{algorithmic}
\end{algorithm}

The detection threshold is determined using the statistical formula given in Equation \ref{eq:threshold}:

\begin{equation}
T = \mu + 3\sigma \quad \text{(Equation 4.6)}
\label{eq:threshold}
\end{equation}

Equation 4.6 denotes the anomaly detection threshold computed as the mean plus three standard deviations from the benign training distribution, ensuring that approximately 99.7\% of normal behavior falls below this boundary and only significant deviations are flagged as anomalies. where $\mu$ is the mean anomaly score computed from benign graphs and $\sigma$ is the standard deviation. The system uses a fixed threshold parameter $k = 3$, corresponding to 3 standard deviations from the mean. This provides a statistically sound basis for anomaly detection, following the 68-95-99.7 rule where approximately 99.7\% of benign behavior falls below the threshold. This approach allows the threshold to adapt to the specific characteristics of the monitored environment while maintaining a principled decision boundary.



\section{\uppercase{ALERT GENERATION AND LOGGING}}

When an anomaly is detected, the system generates structured alerts that provide actionable information to system administrators and security analysts. Alert generation integrates detection results with contextual system information to facilitate rapid incident response.

The alert generation and logging mechanism is presented in Algorithm 4.6.

\begin{algorithm}
\caption{Alert Generation and Logging}
\begin{algorithmic}[1]

\State Input detection result, anomaly score, graph structure

\State Create alert object with timestamp
\State Assign severity level based on anomaly score

\If{anomaly score > critical threshold}
\State Set priority = HIGH
\Else
\State Set priority = MEDIUM
\EndIf

\State Extract associated process information
\State Extract system resource information
\State Compile affected entities list

\For{each graph edge}
\State Record involved processes and resources
\EndFor

\State Serialize alert to structured format (JSON)
\State Log alert to persistent storage
\State Push alert to real-time notification queue

\State Return alert record

\end{algorithmic}
\end{algorithm}

Alerts are logged in JSON format for compatibility with external security information and event management (SIEM) systems. The logging system maintains event history regardless of alert severity, enabling forensic analysis and pattern identification over extended time periods.

Priority levels are dynamically assigned based on anomaly scores rather than using fixed categories, ensuring that alerts accurately reflect the confidence in detection results. This approach minimizes alert fatigue by distinguishing between high-confidence detections and tentative anomalies.



\section{\uppercase{WEB DASHBOARD AND VISUALIZATION}}

The system provides a web-based dashboard that enables security analysts to interact with detection results and monitor system behavior in real time. The dashboard is implemented using Flask and displays behavioral graphs, system metrics, and alert history.

The dashboard architecture and visualization pipeline are described in Algorithm 4.7.

\begin{algorithm}
\caption{Web Dashboard Visualization and Interaction}
\begin{algorithmic}[1]

\State Initialize Flask web server on port 5000
\State Connect to data backend

\While{server is running}

\State Wait for client request
\State Parse request type (graph, metrics, alerts)

\If{request type = graph visualization}
\State Retrieve requested graph
\State Apply layout algorithm
\State Render graph as interactive visualization
\State Send to client

\Else If{request type = system metrics}
\State Compute current system statistics
\State Format metrics for display
\State Send time-series data
\Else If{request type = alert history}
\State Query alert log
\State Filter by date range and severity
\State Send sorted alert list
\EndIf

\EndWhile

\end{algorithmic}
\end{algorithm}

The visualization component uses NetworkX for graph layout algorithms and Matplotlib for rendering. Interactive features enable analysts to inspect individual graph structures and identify suspicious node clusters, filter alerts by severity timestamp and affected resources, correlate detection results with system metrics, export detection results for external analysis, and configure detection parameters and thresholds. These capabilities provide comprehensive interaction with the detection framework, enabling users to explore behavioral patterns and customize system behavior to their operational requirements.

The dashboard updates in real time as new detection results become available, providing immediate visibility into emerging security threats. Historical data enables trend analysis and identification of recurring anomalies.



\section{\uppercase{CONTINUAL LEARNING AND MODEL UPDATES}}

To maintain effectiveness as system behavior evolves, the proposed system incorporates continual learning mechanisms that incrementally adapt the detection model with newly observed benign behavior.

The continual learning procedure is detailed in Algorithm 4.8.

\begin{algorithm}
\caption{Continual Learning and Model Adaptation}
\begin{algorithmic}[1]

\State Input streaming behavior graphs $G_{stream}$
\State Load current model $M_{current}$
\State Initialize replay memory buffer (capacity: 64 graphs)

\State Define adaptation chunk size $C = 3$ graphs
\State Define replay ratio $\rho = 0.5$
\State Define inner epochs per chunk: 2

\For{each chunk of $C$ graphs}

\State Sample replay graphs: $G_{replay} \sim \rho \cdot |C|$ from memory
\State Combine current chunk with replayed graphs

\For{inner epoch = 1 to 2}
\State Fine-tune model on combined batch
\State Compute chunk loss
\EndFor

\State Update memory buffer with current chunk
\State If memory exceeds capacity:
\State Remove oldest graphs

\EndFor

\State Log adaptation statistics

\end{algorithmic}
\end{algorithm}
\newpage
Continual learning allows the model to adapt to gradual shifts in system behavior by processing incoming graphs in small chunks of 3 graphs. This approach reduces concept drift where system behavior patterns gradually change over time. Memory-based replay prevents catastrophic forgetting by interleaving historical benign examples with new data, while remaining conservative to prevent poisoning attacks where attackers inject malicious patterns.



\section{\uppercase{SYSTEM PERFORMANCE OPTIMIZATION}}

Achieving real-time anomaly detection with minimal latency and resource consumption requires careful system optimization. The implementation employs several performance enhancements to reduce computational overhead while maintaining detection accuracy.

The performance optimization strategy is presented in Algorithm 4.9.

\begin{algorithm}
\caption{Performance Optimization and Latency Management}
\begin{algorithmic}[1] % Added [1] for line numbering consistency
\State Initialize system with default buffer sizes
\State Start monitoring performance metrics
\While{system is active}
    \State Measure event collection latency
    \State Measure graph construction time
    \State Measure model inference time
    \State Measure total pipeline latency
    \State Check CPU and memory utilization
    \If{total latency > 1000 ms}
        \State Enable buffer batching
        \State Reduce model update frequency
    \elself{CPU utilization > 80}
        \State Reduce polling frequency
        \State Implement event sampling
    \elself{memory usage > 85}
        \State Shrink event buffer capacity
        \State Implement graph caching strategy
    \EndIf
\EndWhile
\end{algorithmic}
\end{algorithm}

Key optimization techniques implemented in the system include event buffering, which collects events in batches to reduce context switching overhead; graph caching, which maintains an in-memory cache of recent graphs to avoid reconstruction overhead; sparse matrix representation, which utilizes sparse adjacency matrices for memory efficiency; GPU acceleration, which offloads graph autoencoder computations to GPU when available to reduce inference time; and incremental updates, which compute graph updates incrementally rather than reconstructing entire graphs from scratch. These techniques work synergistically to maintain the performance characteristics required for real-time detection.

The system achieves average end-to-end detection latency of 200-400 milliseconds, enabling real-time alerting for emerging threats.



\section{\uppercase{VALIDATION METHODOLOGY}}

The effectiveness of the proposed intrusion detection system is rigorously evaluated through controlled experiments using both simulated and recorded attack scenarios. The validation methodology ensures that detection performance is systematically measured across diverse threat types.

The validation and testing procedure is described in Algorithm 4.10.

\begin{algorithm}
\caption{System Validation and Testing Methodology}
\begin{algorithmic}[1]

\State Collect benign baseline events
\State Build benign training graphs
\State Train autoencoder model

\State Generate simulated attack scenarios

\For{each attack type}

\State Execute attack simulation
\State Collect generated events
\State Construct behavior graphs
\State Apply detection model
\State Record detection results

\State Compute true positive rate (TPR)
\State Compute false positive rate (FPR)
\State Compute detection latency
\State Compute F1-score

\EndFor

\State Aggregate performance metrics
\State Generate ROC and precision-recall curves
\State Analyze false positive patterns
\State Document detection effectiveness

\State Return evaluation report

\end{algorithmic}
\end{algorithm}

The validation procedure evaluates the system against multiple attack types including reverse shells, privilege escalation attempts, data exfiltration, and process injection attacks. Separate benign datasets ensure that performance metrics accurately reflect detection capability on unseen data.

Performance is measured using standard metrics: True Positive Rate (TPR) represents the proportion of attacks correctly detected by the system; False Positive Rate (FPR) represents the proportion of benign instances incorrectly flagged as anomalies; F1-Score measures the harmonic mean of precision and recall, balancing detection accuracy against false alarm rate; and Detection Latency measures the time from attack initiation to alert generation, ensuring the system responds in real time to threats. These metrics provide comprehensive evaluation of both detection effectiveness and operational suitability.



\section{\uppercase{ATTACK SIMULATION AND EVALUATION}}

To evaluate the effectiveness of the proposed intrusion detection system, simulated attack scenarios are generated and executed in a controlled environment. The attack simulation framework provides a safe method to test system capabilities without exposing production systems to real threats.

The evaluation procedure is summarized in Algorithm 4.11.

\begin{algorithm}
\caption{Attack Simulation and Evaluation}
\begin{algorithmic}[1]

\State Generate simulated attack scenario

\If{attack type is Reverse Shell}
\State Spawn shell process
\State Establish outbound network connection
\EndIf

\If{attack type is Privilege Escalation}
\State Access protected files
\State Attempt permission modification
\EndIf

\If{attack type is Data Exfiltration}
\State Read large number of files
\State Send data to external server
\EndIf

\If{attack type is Process Injection}
\State Inject malicious code into running process
\State Execute injected code
\EndIf

\State Record generated events
\State Construct behavior graphs
\State Apply anomaly detection model

\State Display detection results

\end{algorithmic}
\end{algorithm}

The evaluation experiments demonstrate the capability of the proposed system to detect malicious behavior patterns in real time across diverse attack vectors.

\section{\uppercase{IMPLEMENTATION INTEGRATION AND DEPLOYMENT}}

The complete implementation brings together all the proposed components into a unified intrusion detection framework that operates continuously on host systems. The modular architecture enables each component to function independently while maintaining tight integration through standardized data formats and interfaces. Real-time event collection operates in parallel with graph construction and anomaly detection, ensuring minimal processing latency between attack execution and alert generation. The continual learning mechanism adapts the autoencoder model incrementally as new benign behavior patterns emerge, preventing model degradation while maintaining resistance to adversarial data injection. The web-based dashboard provides security analysts with actionable intelligence regarding detected anomalies and system behavior trends, enabling rapid incident response and forensic investigation. The implementation is designed to operate continuously with minimal computational overhead, operating entirely within user space without requiring kernel modifications or elevated system privileges. 




